% theme: quadratic_equations_solid_geometry_applications
\MajorQuestionLesson{立体の利用問題}
\SetRowSpace{6mm}

\TypeQuestionSingle{次の問いに答えなさい。}{chokuhoutai}
%直方体以外(錐体など)の問題も入れてください。
\QQ{底面の横が縦より3cm長い直方体がある。高さが5cm、体積が270cm$^3$であるとき、底面の縦と横の長さを求めなさい。}{縦$6$cm、横$9$cm}
\QQ{底面の縦と横の長さの和が17cmである直方体がある。高さが4cm、体積が240cm$^3$であるとき、底面の縦と横の長さを求めなさい。}{$5$cm、$12$cm}
\QQ{底面の横が縦より4cm長い直方体がある。高さが3cm、体積が168cm$^3$であるとき、底面の縦と横の長さを求めなさい。}{縦$2\sqrt{15}-2$cm、横$2\sqrt{15}+2$cm}

%円錐の問題も入れてください。
\QQ{半径が $x$ cm、高さが6cmの円柱がある。この円柱の体積が150$\pi$ cm$^3$であるとき、$x$ の値を求めなさい。}{$x=5$}
\QQ{ある円柱の高さを変えずに半径だけを2cm長くすると、体積がもとの2倍になった。もとの円柱の半径を求めなさい。}{$2+2\sqrt{2}$cm}
\QQ{ある円柱の高さを変えずに半径だけを1cm短くすると、体積がもとの半分になった。もとの円柱の半径を求めなさい。}{$2+\sqrt{2}$cm}


\TypeQuestionSingle{次の問いに答えなさい。}{kumiawase}
\QQ{1辺$x$ cmの正方形を底面とする高さ4cmの直方体から、1辺2cmの正方形を底面とする高さ4cmの直方体を取り除いたところ、残った立体の体積が80cm$^3$になった。$x$ の値を求めなさい。}{$x=2\sqrt{6}$}
\QQ{1辺の長さが3cm違う2つの立方体がある。2つの立方体の表面積の和が390cm$^2$であるとき、それぞれの立方体の1辺の長さを求めなさい。}{$4$cm、$7$cm}

\TypeQuestionSingle{次の問いに答えなさい。}{youki}
%形式が同じすぎです。
%縦が横より〜〜cmだけ長い(とか短い)タイプの問題
%元の紙が正方形の問題、などを追加して。
%テーマも紙だけじゃないほうが良い。
\QQ{縦と横の長さの比が $2:3$ の長方形の紙がある。この紙の4すみから1辺が3cmの正方形を取り除き、直方体の容器を作ったところ、容積が540cm$^3$になった。この紙の縦と横の長さを求めなさい。}{縦$16$cm、横$24$cm}
\QQ{縦と横の長さの比が $3:4$ の長方形の紙がある。この紙の4すみから1辺が2cmの正方形を取り除き、直方体の容器を作ったところ、容積が560cm$^3$になった。この紙の縦と横の長さを求めなさい。}{縦$18$cm、横$24$cm}


\LessonPageBreak

\TypeQuestionSingle{次の問いに答えなさい。}{sankakuchu}
%test側と全く同じですが、若干変えてください。
%$\Lseg{DP}=\Lseg{EQ}$ではなくて、$\Lseg{PE}=\Lseg{EQ}$にするとか。
\QQSideFig{60mm}{fig/test_07_01}{図のように、底面が直角三角形である直三角柱 ${\rm ABC-DEF}$ があります。$\Ang{ABC}=90^\circ$、$\Lseg{AB}=6$cm、$\Lseg{BC}=8$cm、$\Lseg{AD}=5$cmです。辺 $\Lseg{DE}$ 上に点 ${\rm P}$、辺 $\Lseg{EF}$ 上に点 ${\rm Q}$ をとり、$\Lseg{DP}=\Lseg{EQ}$ とします。三角錐 ${\rm A-DPQ}$ の体積が、この三角柱の体積の $\dfrac{1}{16}$ になるとき、$\Lseg{DP}$ の長さを求めなさい。}{3cm}
\QQ{縦と横の長さの比が $2:3$ の長方形の紙があります。この紙の4すみから1辺が3cmの正方形を切り取り、直方体の容器を作ったところ、容積が540cm$^3$になりました。この紙の縦と横の長さを求めなさい。}{縦16cm、横24cm}
\EndTwoCols


\TypeQuestionSingle{次の問いに答えなさい。}{chokuhoutai}
%test側と同じ感じの問題ですが、直方体を切断し、角の三角錐を切り取る問題に。
%直方体の１つの点と通り、２辺の間を通るような平面で切る。

